% Estilo comum dos esquemáticos do relatório (circuitikz).
% Uso: \documentclass[11pt,border=4pt]{standalone} e, em seguida, % Estilo comum dos esquemáticos do relatório (circuitikz).
% Uso: \documentclass[11pt,border=4pt]{standalone} e, em seguida, % Estilo comum dos esquemáticos do relatório (circuitikz).
% Uso: \documentclass[11pt,border=4pt]{standalone} e, em seguida, % Estilo comum dos esquemáticos do relatório (circuitikz).
% Uso: \documentclass[11pt,border=4pt]{standalone} e, em seguida, \input{estilo_circuitos}.
\usepackage[T1]{fontenc}
\usepackage[utf8]{inputenc}
\usepackage[brazilian]{babel}
\usepackage{lmodern}
\usepackage{amsmath}
\usepackage[american]{circuitikz}
\usetikzlibrary{calc,arrows.meta}

% Cores das figuras originais
\definecolor{anot}{RGB}{0,135,200}      % anotações: valores calculados, setas, resistências vistas
\definecolor{sombra}{RGB}{220,242,253}  % região do transistor nos modelos T e pi
\definecolor{term}{RGB}{62,124,148}     % nomes de terminais destacados (B, C, E)

% Fios com 0,4 pt e corpo dos componentes com 0,8 pt, como nas figuras originais
\ctikzset{bipoles/thickness=2}
.
\usepackage[T1]{fontenc}
\usepackage[utf8]{inputenc}
\usepackage[brazilian]{babel}
\usepackage{lmodern}
\usepackage{amsmath}
\usepackage[american]{circuitikz}
\usetikzlibrary{calc,arrows.meta}

% Cores das figuras originais
\definecolor{anot}{RGB}{0,135,200}      % anotações: valores calculados, setas, resistências vistas
\definecolor{sombra}{RGB}{220,242,253}  % região do transistor nos modelos T e pi
\definecolor{term}{RGB}{62,124,148}     % nomes de terminais destacados (B, C, E)

% Fios com 0,4 pt e corpo dos componentes com 0,8 pt, como nas figuras originais
\ctikzset{bipoles/thickness=2}
.
\usepackage[T1]{fontenc}
\usepackage[utf8]{inputenc}
\usepackage[brazilian]{babel}
\usepackage{lmodern}
\usepackage{amsmath}
\usepackage[american]{circuitikz}
\usetikzlibrary{calc,arrows.meta}

% Cores das figuras originais
\definecolor{anot}{RGB}{0,135,200}      % anotações: valores calculados, setas, resistências vistas
\definecolor{sombra}{RGB}{220,242,253}  % região do transistor nos modelos T e pi
\definecolor{term}{RGB}{62,124,148}     % nomes de terminais destacados (B, C, E)

% Fios com 0,4 pt e corpo dos componentes com 0,8 pt, como nas figuras originais
\ctikzset{bipoles/thickness=2}
.
\usepackage[T1]{fontenc}
\usepackage[utf8]{inputenc}
\usepackage[brazilian]{babel}
\usepackage{lmodern}
\usepackage{amsmath}
\usepackage[american]{circuitikz}
\usetikzlibrary{calc,arrows.meta}

% Cores das figuras originais
\definecolor{anot}{RGB}{0,135,200}      % anotações: valores calculados, setas, resistências vistas
\definecolor{sombra}{RGB}{220,242,253}  % região do transistor nos modelos T e pi
\definecolor{term}{RGB}{62,124,148}     % nomes de terminais destacados (B, C, E)

% Fios com 0,4 pt e corpo dos componentes com 0,8 pt, como nas figuras originais
\ctikzset{bipoles/thickness=2}
